% ─── Актуальность, цель и задачи ──────────────────────────────────────────
\begin{frame}[shrink=10]{Актуальность, цель и задачи}

  \begin{tcolorbox}[formulabox, title={\faExclamationCircle\ Актуальность}]
    \small Прогнозирование валютных курсов --- ключевая задача для банков, инвесторов и госуправления. Финансовые ряды \textbf{нестационарны}, \textbf{нелинейны} и содержат высокий уровень шума. В литературе \textit{нет консенсуса}: превосходят ли нейросети классические модели.
  \end{tcolorbox}

  \vspace{0.2cm}
  \begin{columns}[T]
    \column{0.48\textwidth}
    \textbf{\color{primary}Цель работы:}
    \begin{tcolorbox}[infobox]
      \small Сравнить эффективность классических статистических моделей и нейросетевых методов при прогнозировании курса USD/UZS.
    \end{tcolorbox}

    \column{0.48\textwidth}
    \textbf{\color{primary}Задачи:}
    \begin{itemize}\small
      \setlength\itemsep{2pt}
      \item Изучить теорию временных рядов
      \item Реализовать ARIMA, Random Forest, LSTM
      \item Провести эксперименты на данных ЦБ РУз
      \item Сравнить модели по MAE, RMSE, MAPE
    \end{itemize}
  \end{columns}
\end{frame}
